% source: manti.01.ma, 1ère SM, série suites (Drive 0B6pfSJcEVng0U2ctTDRsUFdyRFk)
\ex{Suites}{U2ct}{Ex. 1}{12}
Soit $(U_n)_n$ la suite définie par $U_0=\dfrac32$ et $U_{n+1}=\dfrac{2}{3-U_n}$.
\begin{enumerate}
\item Montrer que $(\forall n\in\mathbb{N})\ 1<U_n<2$.
\item Vérifier que $U_{n+1}-U_n=\dfrac{(U_n-1)(U_n-2)}{3-U_n}$ et étudier la monotonie de la suite $(U_n)_n$.
\item On pose $V_n=\dfrac{2-U_n}{1-U_n}$ pour tout $n\in\mathbb{N}$.
\begin{enumerate}
\item Montrer que $(V_n)_n$ est une suite géométrique et calculer $V_n$ en fonction de $n$.
\item En déduire le terme général $U_n$ en fonction de $n$.
\end{enumerate}
\item Pour tout entier naturel non nul $n$, on pose $S_n=\dfrac1n\displaystyle\sum_{p=0}^{n}U_p$.
\begin{enumerate}
\item Montrer que $(\forall n\in\mathbb{N})\ U_{n+1}-1\leq\dfrac23(U_n-1)$.
\item Montrer que $(\forall n\in\mathbb{N})\ U_n-1\leq\dfrac12\left(\dfrac23\right)^n$.
\item En déduire que $(\forall n\in\mathbb{N}^*)\ 1+\dfrac1n\leq S_n\leq1+\dfrac{5}{2n}-\dfrac1n\left(\dfrac23\right)^n$.
\end{enumerate}
\end{enumerate}
\fin
\ex{Suites}{U2ct}{Ex. 2}{10}
Soit $(U_n)_n$ la suite définie par $U_0=0$ et $U_{n+1}=\dfrac{U_n-1}{U_n+3}$.
\begin{enumerate}
\item Montrer que $(\forall n\in\mathbb{N})\ -1<U_n\leq0$.
\item Étudier la monotonie de la suite $(U_n)_n$.
\item On pose $V_n=\dfrac{1}{1+U_n}$ pour tout $n\in\mathbb{N}$.
\begin{enumerate}
\item Montrer que $(V_n)_n$ est une suite arithmétique et calculer $V_n$ en fonction de $n$.
\item Déterminer le terme général $U_n$ en fonction de $n$.
\end{enumerate}
\item Pour tout entier naturel non nul $n$, on pose $S_n=\dfrac{1}{n^2}\displaystyle\sum_{k=0}^{n}V_k$. Montrer que $(\forall n\in\mathbb{N}^*)\ \left|S_n-\dfrac14\right|\leq\dfrac3n$.
\end{enumerate}
\fin
\ex{Suites}{U2ct}{Ex. 3}{9}
On considère la suite $(U_n)_n$ définie pour tout entier naturel non nul $n$ par $U_n=\displaystyle\sum_{k=1}^{n}\frac{k}{3^k}$.
\begin{enumerate}
\item Calculer $U_1$ et $U_2$.
\item \begin{enumerate}
\item Montrer par récurrence que $(\forall p\geq2)\ p\leq\left(\dfrac32\right)^p$.
\item En déduire que la suite $(U_n)_n$ est majorée par $1$.
\end{enumerate}
\end{enumerate}
\fin
